\section{Designing a Republic}

In our time, no one considers the acquisition or conquest of a land as a basic right. Rather, it must be justified (i.e., made just) in some way. It is instructive to follow news accounts and to discern what is common in all such justifications.

Any thinker considering the creation of a new nation, republic, or imperium needs to start with Plato's \emph{Republic}. It is a thought experiment in which Socrates delineates the characteristics of the ideal state, that is, the state in which perfect justice reigns. The rulers and citizens are not concerned about that, since they take the state for granted, as a given. For them, justice consists of creating and following the laws of the state.

For the Philosopher, however, justice involves the right order of society. This order is natural or organic, unlike the positive law which is conventional (i.e., the result of human conventions). Specifically, this means ``minding one's own business", that is, everyone performs his own proper role in the state. He should expect the ``just desserts" of his efforts. Unfortunately, if the citizens are not virtuous, they will not mind their own business and they will become envious of others, no longer satisfied with their natural role. Since men fall naturally into three distinct types, depending on the dominant center of their soul life, the virtue required for each type is different, as shown below:

\begin{itemize}
\item \textbf{Sensations}. The virtue necessary for those ruled by the appetitive part of the soul is \textbf{temperance}. 
\item \textbf{Emotions}. The virtue necessary for those ruled by the emotional part of the soul is \textbf{courage}. 
\item \textbf{Intelligence}. The virtue necessary for those ruled by the rational part of the soul is \textbf{wisdom}. 
\end{itemize}
\textbf{Justice}, i.e., true objectivity and truthfulness, is possible only for the integrated man who manages to balance the three parts of the soul into a harmonious whole.

\paragraph{Artifacts and Essences}
First it is necessary to grasp the distinction between real essences and artifacts. The former constitute the organic and the natural, the latter, the conventional and artificial.

A real essence, subsisting in the mind of God, represents a unified being, one whose unity is in itself. Hence a squirrel, or a particular subspecies, is an essence, as is a human being. Recalling what Garrigou-Lagrange wrote about angelic intelligences, essences are further organized into larger systems or wholes. Hence, cities, nations, and so on, are also essences (that is why nations were said to be ruled by a specific archangel), although the modern mind tends to see those groupings as accidental, rather than essential. This seems to be true even of explicitly Thomist philosophers.

An artifact, on the other hand, is a human creation and exists in a human mind. An artifact does not possess a unity-in-itself. Rather, it is given a unity by its creator. You will often see simplistic examples of the type describing, say, the ``idea of a table", with tedious discussions on what exactly a table is. Obviously, this is not what is meant, since a table is an artifact not a natural kind.

To be complete, it should be pointed out that qualia also do not represent essences. Thus, there is no ``idea of red", which is known by the senses, not the intellect. This is a simple guide:

\begin{itemize}
\item \textbf{Essences} are known by the higher, intuitive, part of the mind. This is gnosis, a direct knowing. 
\item \textbf{Artifacts} are known by the rational, thinking, part of the mind. They are either created or defined. 
\item \textbf{Qualia} are known directly by the senses and have no part in this discussion. 
\end{itemize}
\paragraph{The Philosopher's Dilemma}
To be clear, we are using the word ``Philosopher" in the specific sense of the man or woman who transcends the confluence of ordinary, transient life, and is able to contemplate the timeless essences. In thinking the ideal state, he runs into an immediate dilemma. All the successful states (in Plato's time) are natural and organic, while the state envisioned or defined by the philosopher is at best an artifact. For Socrates, the just state would require the elimination of natural things such as marriage, familial affections, and the right to property. This cannot be achieved in thought alone, but requires a myth.

Now all the Greek city-states had a myth of their founding. A mythical figure created its constitution which justified the organic structure of the state. These often involved magical stories, such as the original constitution being hidden in a tree. These various myths cannot all be true, at least in a way that would satisfy a philosopher; as we defined it several years ago\footnote{\url{http://www.gornahoor.net/?p=50}}: \emph{A myth is what other people believe}.

Nevertheless, Socrates was forced to resort to a myth as necessary for the ideal state. He called it the ``noble lie". This must not be understood cynically as recommending that the rulers deliberately lie to the citizens to achieve their ends, for these reasons.

\begin{itemize}
\item There is the obvious objection that not every lie is a noble lie. 
\item The rulers are never philosophers so the liar is not noble 
\item The aim of the noble lie is justice, whereas the objective of the ruler's lie is never just 
\end{itemize}
Moreover, it is a strange sort of lie, since the rulers themselves must believe it.

\paragraph{Interlude}
The next and concluding part of this essay will clarify these points:

\begin{itemize}
\item Why lies are necessary in the world 
\item What exactly is a ``myth" 
\item The elements of the noble lie 
\item Why the Philosopher cannot create the ideal state 
\end{itemize}


\flrightit{Posted on 2014-12-04 by Cologero }

\begin{center}* * *\end{center}

\begin{footnotesize}\begin{sffamily}



\texttt{David on 2014-12-04 at 00:48 said: }

Very interesting. Looking foward to the next parts. Especially that last part. Is there any text in the meantime you would suggest reading on the subject ?


\hfill

\texttt{Matt on 2014-12-04 at 14:51 said: }

David,

Only speaking for myself, a possible suggestion would be The City and Man by Leo Strauss. Don't be put off by the accusations thrown at Strauss, or the liking many neoconservatives have for Strauss; he is probably one of the most misunderstood philosophers – both by his opponents and his neocon supporters – of the modern era.

That text should hopefully help to get you started on this subject.


\hfill


\end{sffamily}\end{footnotesize}
